% !PRINTED 367
\weberpage{367}
Ist $D$ positiv, so ist noch $(E, E')$ zu bestimmen.

Es sei
\[
\varepsilon = \frac{t + u\sqrt{D}}{2}
\]
die fundamentale Einheit in $O$ und $\lambda$ der kleinste positive Exponent, f\"ur den
\[\varepsilon^{\lambda} = \frac{t + u\sqrt{D}}{2}\]
% [sic] Die Quelle lässt diese zweite Darstellung unnummeriert, obwohl der folgende Satz auf (17) verweist, und verwendet darin erneut t,u; beides bleibt erhalten.
in $O'$ enthalten ist. Dann besteht $E$ aus den Potenzen $\pm \varepsilon^{\nu}$, und $E'$ aus den Potenzen $\pm \varepsilon^{\lambda\nu}$, und es ist
\[
(E, E') = \lambda,
\]
und die Formel (15) gibt die Klassenzahl $h'$. Eine weitere Bestimmung l\"a\ss t sich im allgemeinen f\"ur $\lambda$ nicht geben.

Ist wieder $\Delta \equiv 1 \, (\mathrm{mod}\ 4)$ und $Q = 2$, so ist $\lambda = 1$ oder $= 3$; n\"amlich $= 1$, wenn in (17) $u$ gerade ist, und $= 3$, wenn $u$ ungerade ist. Letzteres kann nur vorkommen, wenn $\Delta \equiv 5 \, (\mathrm{mod}\ 8)$ ist, und folglich ist f\"ur $\Delta \equiv 1 \, (\mathrm{mod}\ 8)$ immer $\lambda = 1$, w\"ahrend f\"ur $\Delta \equiv 1 \, (\mathrm{mod}\ 8)$ % [sic]
$\lambda$ sowohl $= 1$ als $= 3$ sein kann.

\begin{center}\rule{0.32\linewidth}{0.4pt}\end{center}
